\documentclass[12pt]{article}

\usepackage{amsmath}
\usepackage{graphicx}
\usepackage{booktabs}
\usepackage[a4paper, margin=1in]{geometry}

\title{Labwork 2: Linear Regression}
\author{Nguyen Xuan Vinh}
\date{\today}

\begin{document}

\maketitle

\section{Introduction}

Linear Regression is one of the most fundamental supervised learning algorithms in Machine Learning. 
It is used to model the relationship between an input variable and an output variable using a linear equation.

In this labwork, Linear Regression was implemented completely from scratch using Gradient Descent optimization without using external Machine Learning libraries.

The objective is to predict house prices based on house size using a linear model.

\section{Mathematical Formulation}

The Linear Regression model is defined as:

\[
\hat{y} = w_1x + w_0
\]

where:

\begin{itemize}
    \item \(x\) is the input feature (house size)
    \item \(\hat{y}\) is the predicted house price
    \item \(w_1\) is the weight (slope)
    \item \(w_0\) is the bias (intercept)
\end{itemize}

To optimize the parameters, Mean Squared Error (MSE) was used as the loss function:

\[
MSE = \frac{1}{n}\sum_{i=1}^{n}(y_i - \hat{y}_i)^2
\]

The gradients of the loss function are:

\[
\frac{\partial MSE}{\partial w_1}
=
\frac{-2}{n}\sum x_i(y_i - \hat{y}_i)
\]

\[
\frac{\partial MSE}{\partial w_0}
=
\frac{-2}{n}\sum (y_i - \hat{y}_i)
\]

The Gradient Descent update rules are:

\[
w_1 = w_1 - r \frac{\partial MSE}{\partial w_1}
\]

\[
w_0 = w_0 - r \frac{\partial MSE}{\partial w_0}
\]

where \(r\) is the learning rate.

\section{Algorithm Implementation}

The algorithm was implemented manually in Python.

The implementation process consists of the following steps:

\begin{enumerate}
    \item Read training data from the CSV file \texttt{lr.csv}
    \item Initialize parameters \(w_1\) and \(w_0\)
    \item Compute predictions using the linear equation
    \item Compute prediction errors
    \item Calculate gradients for \(w_1\) and \(w_0\)
    \item Update parameters using Gradient Descent
    \item Repeat the process iteratively
    \item Print intermediate optimization results
\end{enumerate}

The implementation does not use libraries such as Scikit-learn, TensorFlow, or PyTorch.

\section{Dataset}

The dataset contains house sizes and corresponding house prices.

\begin{table}[h]
\centering
\begin{tabular}{cc}
\toprule
House Size & House Price \\
\midrule
10 & 55 \\
20 & 80 \\
40 & 100 \\
60 & 120 \\
80 & 150 \\
\bottomrule
\end{tabular}
\caption{Training dataset from lr.csv}
\end{table}

\section{Experimental Results}

The model was trained using Gradient Descent.

The initial values of the parameters were:

\[
w_1 = 0,\quad w_0 = 0
\]

The learning rate used in the main experiment was:

\[
r = 0.0001
\]

\subsection{Intermediate Optimization Results}

The following table shows several intermediate training steps produced by the program.

\begin{table}[h]
\centering
\begin{tabular}{cccc}
\toprule
Iteration & \(w_1\) & \(w_0\) & Loss \\
\midrule
0 & 1.0140 & 0.0202 & 11265.0000 \\
100 & 2.0900 & 0.2981 & 635.4750 \\
200 & 2.0855 & 0.5559 & 628.8275 \\
300 & 2.0810 & 0.8123 & 622.2516 \\
400 & 2.0766 & 1.0673 & 615.7466 \\
500 & 2.0722 & 1.3209 & 609.3118 \\
600 & 2.0678 & 1.5732 & 602.9463 \\
700 & 2.0635 & 1.8241 & 596.6494 \\
800 & 2.0591 & 2.0736 & 590.4205 \\
900 & 2.0548 & 2.3218 & 584.2587 \\
\bottomrule
\end{tabular}
\caption{Intermediate Gradient Descent optimization results}
\end{table}

The final learned parameters are:

\[
w_1 = 2.0506
\]

\[
w_0 = 2.5661
\]

The loss value continuously decreases during training, showing that Gradient Descent successfully optimizes the Linear Regression model.

\subsection{Effect of Learning Rate}

Different learning rates produce different convergence behaviors.

\begin{itemize}
    \item Small learning rates produce stable but slow convergence
    \item Suitable learning rates improve optimization efficiency
    \item Large learning rates may cause unstable updates and divergence
\end{itemize}

For this experiment, the learning rate \(r = 0.0001\) provides stable convergence and gradual reduction of the loss function.

\section{Discussion}

The experiments demonstrate that Gradient Descent can effectively optimize the parameters of a Linear Regression model.

As training progresses:
\begin{itemize}
    \item The loss value decreases
    \item The predicted values become more accurate
    \item The parameters gradually approach optimal values
\end{itemize}

The learning rate is an important hyperparameter because it directly affects convergence speed and training stability.

\section{Conclusion}

In this labwork, Linear Regression was successfully implemented from scratch using Gradient Descent optimization.

The program:
\begin{itemize}
    \item Reads training data from a CSV file
    \item Predicts house prices using a linear equation
    \item Optimizes parameters iteratively
    \item Prints intermediate optimization results
\end{itemize}

The final experimental results show that Gradient Descent successfully minimizes the prediction error and learns the relationship between house size and house price.

This labwork provides practical understanding of optimization techniques used in Machine Learning and Deep Learning.

\end{document}